\documentclass[10pt,a4paper]{amsart}
\usepackage{amsmath,amssymb,amsthm,graphicx}
% Retained preamble for the Stacks Project chapter divisors.tex (Divisors).
% The chapter file pinned-original-source.tex carries only \input{preamble} on its first line, so no
% preamble text lives inside the chapter itself.  The Stacks Project uses one single shared preamble for
% every chapter; the snapshot stages it verbatim in sources/stacks-chow/preamble.tex.  The block below is
% copied mechanically, line for line, from that file's lines 14-18 (the Xy-pic setup) and lines 154-217
% (the native theorem declarations together with the mathematical macros).  Nothing was added, removed or
% redefined.  The wrapper class and the amsmath/amssymb/amsthm/graphicx base packages are supplied by the
% assembler, so the class/\externaldocument part of the shared preamble is not repeated here.
% For commutative diagrams we use Xy-pic
\usepackage[all]{xy}

% We use 2cell for 2-commutative diagrams.
\xyoption{2cell}
% Theorem environments.
%
\theoremstyle{plain}
\newtheorem{theorem}[subsection]{Theorem}
\newtheorem{proposition}[subsection]{Proposition}
\newtheorem{lemma}[subsection]{Lemma}

\theoremstyle{definition}
\newtheorem{definition}[subsection]{Definition}
\newtheorem{example}[subsection]{Example}
\newtheorem{exercise}[subsection]{Exercise}
\newtheorem{situation}[subsection]{Situation}

\theoremstyle{remark}
\newtheorem{remark}[subsection]{Remark}
\newtheorem{remarks}[subsection]{Remarks}

\numberwithin{equation}{subsection}

% Macros
%
\def\lim{\mathop{\mathrm{lim}}\nolimits}
\def\colim{\mathop{\mathrm{colim}}\nolimits}
\def\Spec{\mathop{\mathrm{Spec}}}
\def\Hom{\mathop{\mathrm{Hom}}\nolimits}
\def\Ext{\mathop{\mathrm{Ext}}\nolimits}
\def\SheafHom{\mathop{\mathcal{H}\!\mathit{om}}\nolimits}
\def\SheafExt{\mathop{\mathcal{E}\!\mathit{xt}}\nolimits}
\def\Sch{\mathit{Sch}}
\def\Mor{\mathop{\mathrm{Mor}}\nolimits}
\def\Ob{\mathop{\mathrm{Ob}}\nolimits}
\def\Sh{\mathop{\mathit{Sh}}\nolimits}
\def\NL{\mathop{N\!L}\nolimits}
\def\CH{\mathop{\mathrm{CH}}\nolimits}
\def\proetale{{pro\text{-}\acute{e}tale}}
\def\etale{{\acute{e}tale}}
\def\QCoh{\mathit{QCoh}}
\def\Ker{\mathop{\mathrm{Ker}}}
\def\Im{\mathop{\mathrm{Im}}}
\def\Coker{\mathop{\mathrm{Coker}}}
\def\Coim{\mathop{\mathrm{Coim}}}

% Boxtimes
%
\DeclareMathSymbol{\boxtimes}{\mathbin}{AMSa}{"02}

%
% Macros for moduli stacks/spaces
%
\def\QCohstack{\mathcal{QC}\!\mathit{oh}}
\def\Cohstack{\mathcal{C}\!\mathit{oh}}
\def\Spacesstack{\mathcal{S}\!\mathit{paces}}
\def\Quotfunctor{\mathrm{Quot}}
\def\Hilbfunctor{\mathrm{Hilb}}
\def\Curvesstack{\mathcal{C}\!\mathit{urves}}
\def\Polarizedstack{\mathcal{P}\!\mathit{olarized}}
\def\Complexesstack{\mathcal{C}\!\mathit{omplexes}}
% \Pic is the operator that assigns to X its picard group, usage \Pic(X)
% \Picardstack_{X/B} denotes the Picard stack of X over B
% \Picardfunctor_{X/B} denotes the Picard functor of X over B
\def\Pic{\mathop{\mathrm{Pic}}\nolimits}
\def\Picardstack{\mathcal{P}\!\mathit{ic}}
\def\Picardfunctor{\mathrm{Pic}}
\def\Deformationcategory{\mathcal{D}\!\mathit{ef}}

% --- appended 2026-10-12 (Chief): the project own definitions of the reference/slogan/history
% display environments, copied mechanically from the shared Stacks preamble lines 8-12
% (\usepackage{verbatim} plus the three \newenvironment lines, which the chapter preamble keeps
% verbatim from the project preamble). Nothing added, removed or redefined: these environments
% carry no mathematical content and the project itself typesets them as comments. Without these
% three lines a counted statement whose own text uses them fails with Environment slogan undefined.
\usepackage{verbatim}
\newenvironment{reference}{\comment}{\endcomment}
\newenvironment{slogan}{\comment}{\endcomment}
\newenvironment{history}{\comment}{\endcomment}
\usepackage{cleveref}
\crefname{originaldefinition}{original definition}{original definitions}
\Crefname{originaldefinition}{Original definition}{Original definitions}
\makeatletter
\newlabel{algebra-definition-relative-assassin}{{05AT}{source}{}{sourcepointer-1}{}}
\newlabel{algebra-definition-relative-assassin@cref}{{[originaldefinition][0][]05AT}{[1][0][]source}{}{}{}}
\makeatother
\crefname{originallemma}{original lemma}{original lemmas}
\Crefname{originallemma}{Original lemma}{Original lemmas}
\makeatletter
\newlabel{algebra-lemma-compare-relative-assassins}{{05GB}{source}{}{sourcepointer-2}{}}
\newlabel{algebra-lemma-compare-relative-assassins@cref}{{[originallemma][0][]05GB}{[1][0][]source}{}{}{}}
\makeatother
\crefname{originallemma}{original lemma}{original lemmas}
\Crefname{originallemma}{Original lemma}{Original lemmas}
\makeatletter
\newlabel{lemma-associated-affine-open}{{02OK}{source}{}{sourcepointer-3}{}}
\newlabel{lemma-associated-affine-open@cref}{{[originallemma][0][]02OK}{[1][0][]source}{}{}{}}
\makeatother
\title{Relative associated primes over an affine chart detect relative associated points of the sheaf}
\author{The Stacks Project Authors}
\date{stacks-project:divisors.tex:local-snapshot-20261012; GNU FDL 1.2 or later}
\begin{document}
\maketitle
\noindent\textit{Editorial note.} All delivered mathematical text is the author's own, copied from the pinned Stacks Project chapter file pinned-original-source.tex; no statement or proof step was rewritten, reordered or removed.\par\medskip
\noindent\textit{Editorial note.} Following the multi-claim convention for this source, source\_url is the source-level URL https://stacks.math.columbia.edu/ and the source label lemma-relative-assassin-affine-open is carried in claim\_key and here rather than in a per-claim tag URL.\par\medskip
\noindent\textit{Editorial note.} The preamble is the chapter's shared Stacks Project preamble, retained as sources/stacks-divisors/retained-preamble.tex: a line-for-line excerpt of the shared Stacks preamble staged with this snapshot. No macro was invented or redefined, and this card adds no preamble text of its own.\par\medskip
\noindent\textit{Editorial note.} Statement span 1078-1095 and proof span 1097-1117 come from the halfA pool card operations/pool-cards/stacks-divisors-halfA.json and were re-read in the pinned chapter file before this card was written.\par\medskip
\noindent\textit{Editorial note.} The proof refers to the relative associated primes by the author's notation Ass\_\{A/R\}(M) and to the sheaf version by Ass\_\{X/S\}(F); both are preserved.\par\medskip
\noindent\textit{Editorial note.} The statement's final sentence is an equivalence whose converse needs finite generation of the fibre prime; the proof marks that openly, and the hypothesis 'all fibres locally Noetherian' in the statement is what supplies it.\par\medskip
\noindent\textit{Editorial note.} 3 cross-references occur inside the counted spans; they are listed in external\_references as navigation pointers only.\par\medskip
\noindent\textit{Allowed standard prerequisites (editorial).}\begin{itemize}
\item The definition of the relative associated primes of a module over a ring map, cited from the algebra chapter. Cited at proof lines 1097-1098. Names the algebraic set that the statement compares with the relative associated points of the sheaf.
\item The comparison between relative associated primes of a module and associated primes of its base change to the residue field, cited from the algebra chapter. Cited at proof lines 1105-1107. Turns the algebraic condition into a condition on the fibre of the module, which is where the geometric statement lives.
\item The same-chapter description of associated points of a quasi-coherent sheaf on an affine open by the associated primes of the module of sections. Cited at proof lines 1111-1112, this being the lane's card L4023 of this source. Converts the fibre condition into the statement that the point is associated to the restriction of the sheaf to the fibre.
\end{itemize}\medskip
\ifdefined\hypertarget\hypertarget{sourcepointer-1}{}\fi
\noindent\textit{Source reference.} Original definition 05AT: focused-seeds/AGP/sources/stacks-snapshot-20261011/tags.txt, source lines 6861--6861. Same-chapter reference used in the retained author text at source line(s) 1052, 1099. It is supplied only as a navigation pointer so that the retained text compiles; the counted item delivers the author's own statement and the author's complete proof, and no part of the referenced item is copied in or substituted for the counted argument.\par\medskip
\ifdefined\hypertarget\hypertarget{sourcepointer-2}{}\fi
\noindent\textit{Source reference.} Original lemma 05GB: focused-seeds/AGP/sources/stacks-snapshot-20261011/tags.txt, source lines 7055--7055. Same-chapter reference used in the retained author text at source line(s) 1109. It is supplied only as a navigation pointer so that the retained text compiles; the counted item delivers the author's own statement and the author's complete proof, and no part of the referenced item is copied in or substituted for the counted argument.\par\medskip
\ifdefined\hypertarget\hypertarget{sourcepointer-3}{}\fi
\noindent\textit{Source reference.} Original lemma 02OK: focused-seeds/AGP/sources/stacks-snapshot-20261011/tags.txt, source lines 3519--3519. Same-chapter reference used in the retained author text at source line(s) 168, 186, 423, 852, 1114. It is supplied only as a navigation pointer so that the retained text compiles; the counted item delivers the author's own statement and the author's complete proof, and no part of the referenced item is copied in or substituted for the counted argument.\par\medskip
\begin{lemma}
\label{lemma-relative-assassin-affine-open}
Let $f : X \to S$ be a morphism of schemes.
Let $\mathcal{F}$ be a quasi-coherent sheaf on $X$.
Let $U \subset X$ and $V \subset S$ be affine opens
with $f(U) \subset V$. Write $U = \Spec(A)$, $V = \Spec(R)$, and set
$M = \Gamma(U, \mathcal{F})$.
Let $x \in U$, and let $\mathfrak p \subset A$ be the corresponding prime.
Then
$$
\mathfrak p \in \text{Ass}_{A/R}(M) \Rightarrow
x \in \text{Ass}_{X/S}(\mathcal{F})
$$
If all fibres $X_s$ of $f$ are locally Noetherian, then
$\mathfrak p \in \text{Ass}_{A/R}(M) \Leftrightarrow
x \in \text{Ass}_{X/S}(\mathcal{F})$
for all pairs $(\mathfrak p, x)$ as above.
\end{lemma}

\begin{proof}
The set $\text{Ass}_{A/R}(M)$ is defined in
Algebra, Definition \ref{algebra-definition-relative-assassin}.
Choose a pair $(\mathfrak p, x)$. Let $s = f(x)$.
Let $\mathfrak r \subset R$ be the prime lying under $\mathfrak p$,
i.e., the prime corresponding to $s$.
Let $\mathfrak p' \subset A \otimes_R \kappa(\mathfrak r)$
be the prime whose inverse image is $\mathfrak p$, i.e.,
the prime corresponding to $x$ viewed as a point of its fibre $X_s$.
Then $\mathfrak p \in \text{Ass}_{A/R}(M)$ if and only if
$\mathfrak p'$ is an associated prime of
$M \otimes_R \kappa(\mathfrak r)$, see
Algebra, Lemma \ref{algebra-lemma-compare-relative-assassins}.
Note that the ring $A \otimes_R \kappa(\mathfrak r)$ corresponds to $U_s$
and the module $M \otimes_R \kappa(\mathfrak r)$ corresponds to the
quasi-coherent sheaf $\mathcal{F}_s|_{U_s}$.
Hence $x$ is an associated point of $\mathcal{F}_s$
by Lemma \ref{lemma-associated-affine-open}.
The reverse implication holds if $\mathfrak p'$ is finitely generated
which is how the last sentence is seen to be true.
\end{proof}

\end{document}
